\documentclass[12pt]{article}

% Idioma e codificação
\usepackage[brazil]{babel}
\usepackage[utf8]{inputenc}
\usepackage[T1]{fontenc}

% Matemática e listas
\usepackage{amsmath, amssymb}
\usepackage{enumitem}

% Layout e cores
\usepackage{geometry}
\usepackage{graphicx}
\usepackage{xcolor}

% Cores institucionais
\definecolor{maincolor}{RGB}{0, 79, 151}
\definecolor{meuazul}{RGB}{0, 79, 151}

% Cabeçalho e rodapé
\usepackage{fancyhdr}

% Blocos
\usepackage[most]{tcolorbox}
\newtcolorbox{block}[1]{%
  colback=blue!5,
  colframe=maincolor,
  coltitle=white,
  title=#1,
  fonttitle=\bfseries,
  boxrule=1.5pt,
  arc=3mm,
  left=6pt,
  right=6pt,
  top=6pt,
  bottom=6pt
}

\newtcolorbox{resposta}[1]{%
  colback=white,
  colframe=meuazul!30,
  arc=1mm,
  height=#1,
  width=\textwidth,
  boxrule=0.8pt
}

% Margens
\geometry{margin=2.5cm}

% Fancyhdr
\pagestyle{fancy}
\fancyhf{}

\renewcommand{\headrulewidth}{2pt}
\renewcommand{\footrulewidth}{2pt}

\renewcommand{\headrule}{\hbox to\headwidth{%
  \color{meuazul}\leaders\hrule height \headrulewidth\hfill}}

\renewcommand{\footrule}{\hbox to\headwidth{%
  \color{meuazul}\leaders\hrule height \footrulewidth\hfill}}

\setlength{\headheight}{52pt}
\addtolength{\topmargin}{-20pt}

\fancyhead[L]{%
  \includegraphics[height=1.2cm]{logo.png}
}

\fancyhead[R]{%
  \raggedleft
  \textcolor{meuazul}{%
    \small
    \textbf{Lista de Exercícios — Introdução a Microsserviços}\\
    Me.~Sergio Souza Novak\\
    \textit{Engenharia da Qualidade e Confiabilidade}
  }
}

\fancyfoot[C]{\textcolor{meuazul}{\thepage}}

% =========================
\begin{document}
% =========================

\vspace*{0.2cm}

\begin{center}
  {\Large \textbf{Introdução a Microsserviços}}\\[0.2cm]
  \textit{Instruções: Justifique suas respostas de forma clara e objetiva.
  Para questões discursivas, escreva na caixa de resposta abaixo de cada enunciado.}
\end{center}

\vspace{0.4cm}

% -------------------------------------------------------
% QUESTÃO 1
% -------------------------------------------------------
\textbf{1.} Explique o conceito de arquitetura de microsserviços e cite pelo menos três de suas principais características.
\begin{resposta}{4cm}
\end{resposta}

\vspace{0.5cm}

% -------------------------------------------------------
% QUESTÃO 2
% -------------------------------------------------------
\textbf{2.} Quais foram os principais motivos que levaram grandes empresas da indústria a abandonar a arquitetura monolítica em favor dos microsserviços?
\begin{resposta}{4.5cm}
\end{resposta}

\vspace{0.5cm}

% -------------------------------------------------------
% QUESTÃO 3
% -------------------------------------------------------
\textbf{3.} Compare a arquitetura monolítica e a arquitetura de microsserviços em relação aos seguintes critérios: Implantação, Escalabilidade, Stack Tecnológica, Acoplamento e Resiliência.
\begin{resposta}{5.5cm}
\end{resposta}

\vspace{0.5cm}

% -------------------------------------------------------
% QUESTÃO 4
% -------------------------------------------------------
\textbf{4.} Ao dimensionar (escalar) uma aplicação, qual é a limitação fundamental da abordagem de "processo único"? Como os conceitos de simultaneidade e segmentação resolvem esse problema?
\begin{resposta}{4.5cm}
\end{resposta}

\vspace{0.5cm}

% -------------------------------------------------------
% QUESTÃO 5
% -------------------------------------------------------
\textbf{5.} Apesar de suas vantagens, a adoção de microsserviços introduz novas complexidades. Cite e explique os principais desafios técnicos e organizacionais dessa arquitetura.
\begin{resposta}{5cm}
\end{resposta}

\vspace{0.5cm}

% -------------------------------------------------------
% QUESTÃO 6
% -------------------------------------------------------
\textbf{6.} A migração de um sistema monolítico para microsserviços costuma seguir três etapas principais. Descreva o que deve ser feito em cada uma delas.
\begin{resposta}{5cm}
\end{resposta}

\vspace{0.5cm}

% -------------------------------------------------------
% QUESTÃO 7
% -------------------------------------------------------
\textbf{7.} Em relação à estruturação de equipes (Etapa 2 da migração), quais são as duas metodologias principais de divisão de equipes? Quais os riscos associados à sobrecarga dos desenvolvedores?
\begin{resposta}{4.5cm}
\end{resposta}

\vspace{0.5cm}

% -------------------------------------------------------
% QUESTÃO 8
% -------------------------------------------------------
\textbf{8.} Todo microsserviço possui três componentes em sua arquitetura básica: Frontend (API), Backend e Armazenamento. Explique o papel do Frontend nesse contexto e diferencie-o de uma interface de usuário tradicional (UI).
\begin{resposta}{4.5cm}
\end{resposta}

\vspace{0.5cm}

% -------------------------------------------------------
% QUESTÃO 9
% -------------------------------------------------------
\textbf{9.} Na Camada de Comunicação do ecossistema, os microsserviços podem interagir de forma síncrona ou assíncrona. Descreva a diferença entre usar \textbf{HTTP + REST} (Síncrono) e \textbf{Troca de Mensagens / Pub-Sub} (Assíncrono), destacando vantagens e desvantagens de cada um.
\begin{resposta}{5.5cm}
\end{resposta}

\vspace{0.5cm}

% -------------------------------------------------------
% QUESTÃO 10
% -------------------------------------------------------
\textbf{10.} Por que manter diferentes versões de um mesmo microsserviço (versionamento) é considerado um antipadrão (anti-pattern)?
\begin{resposta}{4cm}
\end{resposta}

\vspace{0.5cm}

% -------------------------------------------------------
% QUESTÃO 11
% -------------------------------------------------------
\textbf{11.} O que o "Paradigma da Caixa-Preta" significa para uma equipe responsável por um microsserviço específico dentro da corporação? Qual o nível de liberdade que os desenvolvedores possuem nesse cenário?
\begin{resposta}{4.5cm}
\end{resposta}

\vspace{0.5cm}

% -------------------------------------------------------
% QUESTÃO 12
% -------------------------------------------------------
\textbf{12.} O ecossistema de microsserviços é sustentado por diferentes camadas. Explique a importância de "separar a infraestrutura" do desenvolvimento prático dos microsserviços e qual é o papel central da Engenharia de Infraestrutura nesse processo.
\begin{resposta}{5cm}
\end{resposta}

\vspace{0.8cm}

\begin{center}
  \textit{``Boa atividade! Att. Me. Prof. Sergio Souza Novak.''}
\end{center}

% =========================
\end{document}
% =========================
